\documentclass{article}
\usepackage{amsmath}
\usepackage{cancel}

\title{The Dissipation Function}
\begin{document}
\maketitle

As we showed previously, dAlembert's principle in it's general form is
\begin{equation}
  \frac{d}{dt}\left( \frac{\partial T}{\partial \dot{q}_j} \right)-\frac{\partial T}{\partial q_j} = Q_j
\end{equation}
for generalized forces $Q_j$. When these generalized forces are derived from a
potential such that
\begin{equation}
  Q_j = \sum_i \vec{F}_i\cdot\frac{\partial \vec{r}_i}{\partial q_j} = -\sum_i\nabla_iV\frac{\partial \vec{r}_i}{\partial q_j} = -\frac{\partial V}{\partial q_j}
\end{equation}
then we can simplify the above to the familiar
\begin{equation}
  \frac{d}{dt}\frac{\partial \mathcal{L}}{\partial \dot{q}_j}-\frac{\partial \mathcal{L}}{\partial q_j} = 0 
\end{equation}
suppose that not all of the forces in our system can be derived from a scalar
potential $V$. Then the Euler-Lagrange equations become
\begin{equation}
  \frac{d}{dt}\frac{\partial \mathcal{L}}{\partial \dot{q}_j}-\frac{\partial \mathcal{L}}{\partial q_j} = Q_j
\end{equation}
where $Q_j$ are the generalized forces that have some kind of velocity
dependence. 

Take for example the case of \textbf{friction}. It is easy to imagine cases
where the frictional force is proportional to the velocity of the moving object
so that in one dimension we would expect
\begin{equation}
  F_{f,x} = -k_xv_x
\end{equation}
This equation appears similar to a Hooke's law force except the coordinate
displacement has been replaced with velocity displacement. Keeping that in mind
we can integrate this force against velocity to obtain a ``potential'' that we
can use for our Lagrangian formulation. 
\begin{equation}
  \mathcal{F} = \frac{1}{2} \sum\limits_{\textrm{particles}}k_jv_j^2
\end{equation}
Such a generalized potential is called \textit{Rayleigh's dissipation function}. 

The generalized force would therefore be
\begin{align}
  Q_j &= \sum_i \vec{F}_{f, i}\cdot\frac{\partial \vec{r}_i}{\partial q_j} \\
      &= \sum_i\frac{\partial \mathcal{F}}{\partial \vec{v}_i}\frac{\partial \dot{\vec{r}}_i}{\partial \dot{q}_j} \\
      &= \sum_i\frac{\partial \mathcal{F}}{\partial \dot{\vec{r}}_i}\frac{\partial \dot{\vec{r}}_i}{\partial \dot{q}_j} \\
      &= -\frac{\partial \mathcal{F}}{\partial \dot{q}_j }
\end{align}
One example where this is useful is in fluid dynamics. Imagine a sphere
traveling through some water. Stoke's law describes the frictional force (drag)
on the sphere by
\begin{equation}
  \vec{F} = 6\pi\eta a \vec{v}
\end{equation}
where $a$ is the radius of the sphere and $\eta$ is the dynamic viscosity. This
is an important law because you can us it to determine the viscosity of a fluid
by dropping a sphere and observing the resulting dynamics. It also can be used
to explain why water droplets in fog can remain suspended until they reach a
critical size; gravity and drag compete.  
\end{document}
